\label{sec:conclusion}

This paper has established a mathematical framework for analysing incumbent pairing
under algorithmic redistricting, derived the analytical baseline expectation for three
representative states, and demonstrated empirically that \textsc{bisect} maps fall
within the geographic random baseline while enacted maps fall significantly below it.

The central finding is stark: enacted congressional maps in NC, WI, and TX all produce
pairing rates below the random redistricting baseline (preliminary single-seed binomial
tests: $p < 0.05$ in all three states, assuming independent pairing events), consistent
with mapmakers having deliberately drawn district boundaries to guarantee each sitting
incumbent a home district. The \textsc{bisect} algorithm, operating without incumbency data,
produces pairing rates consistent with the random baseline across all three states.

For redistricting practitioners, these findings have three direct implications.
First, any redistricting plan that achieves a significantly lower pairing rate than
the geographically adjusted baseline has used incumbency data or has subordinated
other redistricting criteria to achieve pairing protection---a choice that may require
justification in a legal proceeding. Second, \textsc{bisect} maps provide a credible
neutral baseline for pairing-rate comparisons: if an algorithmic map consistent with
geographic constraints produces pairing at the random rate, the enacted map's lower
pairing rate is attributable to mapmaker choice, not geographic necessity. Third, the
incumbency protection premium provides a single number that quantifies the degree to
which a plan departs from geographic neutrality on pairing---useful for expert witness
reports and judicial findings of fact.

Future work should extend this analysis to multi-seed ensemble runs to characterise
the full distribution of pairing outcomes under algorithmic redistricting, and to
commission-drawn and court-drawn maps to establish whether neutral human redistricting
produces pairing rates closer to the algorithmic or enacted baseline.
